
\chapter{Bevezetés}

\section{A dokumentáció felépítése}

A Magyar Ispell dokumentáció első része a
felhasználói ismereteket foglalja össze, második része a
fejlesztők számára tartalmaz útmutatást. A függelék első felében
a Hunspell helyesírás-ellenőrző és -javító program, és az
általa használt szótárállományok formátumának kézikönyvlapjait
találjuk. Itt kapott helyet  egy részletes összehasonlítás
is a Hunspell és a magyar
helyesírási szabályzat (illetve a jelenleg legelterjedtebb magyar
helyesírás-ellenőrző program, a MS Office XP-ben megtalálható
Helyes-e?) között, amit végül egy a Hunspell működésének
helyességét vizsgáló leírás követ.

\section{Elérhetőség}

A Magyar Ispell levelezési listája: \url{magyarispell@yahoogroups.com}

A lista archívuma: \url{http://www.yahoogroups.com/group/magyarispell}

A Magyar Ispell honlapja: \url{http://magyarispell.sourceforge.net}

A Szószablya honlapja: \url{http://www.szoszablya.hu}

\section{Köszönetnyilvánítás}

A dokumentáció korábbi változatának elkészítésében, javításában,
kiegészítésében Mayer Gyula segített.

A fejlesztésben sokan részt vettek. Szántó Tamás, Bíró Árpád, Hilzinger Marcel,
Koblinger Egmont, Godó Ferenc, Trón Viktor, Pásztor György, Noll János,
Bencsáth Boldizsár, Fransoa Holop, Goldmann Eleonóra, Halácsy Péter,
Kornai András és sokak segítségét köszönöm!

Anyagi jellegű támogatást kaptam Mátó Pétertől, a
Typo\TeX{} könyvkiadótól (Mayer Gyula segítségével), a SuSE Linuxtól (Koblinger
Egmont és Hilzinger Marcel segítségével), az UHU-Linuxtól. Jelenleg
az állami támogatással működő Szószablya munkacsoport tagjaként a
Magyar Ispell fejlesztésén is dolgozom.

A tesztprogramok és a dokumentáció 2002. évi változatának
elkészítéséhez, valamint a tesztelés kivitelezéséhez a
Typo\TeX{}\ Kft. 300\,000 (+áfa) forinttal járult hozzá,
a Széchenyi Terv keretében az Informatikai Kormánybiztosság SZ-IS-10/3
pályázati számon elnyert támogatásából. 2002 őszétől az
összegből vásárolt számítógéppel készülnek a Magyar Ispell új változatai.

Az IMEDIA Kft. (\url{http://www.imedia.hu}) különösen nagy segítséget
nyújtott a szókincs ellenőrzésében és bővítésében, azzal, hogy
nagy és ellenőrzött szókincset nyújtott át a fejlesztés számára.

Az UHU-Linux Kft. (\url{http://www.uhulinux.hu}) szoftverfejlesztési pályázatának
támogatásával készült el a 0.98-as Magyar
Ispell és ennek honlapja, valamint a Magyar MySpell használatát lehetővé tevő
AbiWord-programkiegészítés (folt). Noll János és Koblinger Egmont munkája révén az
UHU-Linux 1.0-s változatába is bekerült a legfrissebb Magyar Ispell szótárat és
helyesírás-ellenőrzőt tartalmazó OpenOffice.org irodai csomag, illetve
AbiWord szövegszerkesztő.

A Szószablya (\url{http://www.szoszablya.hu}) révén sor kerül a Magyar
Ispell javítására és tovább fejlesztése is, mind szókincsét, mind a magyar nyelv
kezelését tekintve, egy nyitott forráskódú magyar morfológiai elemző, a Hunmorph
elkészítése céljából. A Szószablya keretén belül elvégzett vizsgálatok szerint
a Magyar Ispell 0.99.3-as változata a magyar weboldalakon található (helyes)
szóalakok körülbelül 96\%-át -- beleértve ebbe az összes gyakori
szót -- felismeri.

\section{A dokumentáció}

A dokumentáció forrása \LaTeX\ nyelven készült. A mesterfájlban a
\verb+\usepackage{times}+ utasítás abból a célból szerepel, hogy a
\texttt{pdftex} program segítségével jó minőségű és kis méretű
PDF változatot lehessen előállítani.

A dokumentáció HTML változata a
\begin{verbatim}
latex2html -split 4 magyarispell
\end{verbatim}
parancs segítségével szinte hiba nélkül előállítható.

\section{Felhasználási engedély}

A Magyar Ispell szótár jelen változatát Németh László készítette
a fejlesztés résztvevőinek és támogatóinak segítségével,
és a GNU GPL (General Public License) kettes, vagy későbbi változata
alapján szabadon felhasználható.

A Hunspell program, és a Magyar Myspell függvénykönyvtár
a GNU LGPL (GNU Library GPL) kettes, vagy későbbi változata
alapján szabadon felhasználható.

Ez a dokumentum a Free Software Foundation által kiadott GNU Free
Documentation License 1.2-es vagy újabb változatában foglalt feltételek
szerint másolható, terjeszthető és módosítható; rögzített
fejezet, első borítólapszöveg és hátsó borítólapszöveg nincsen
(http://www.gnu.org/licenses/fdl.html).

Németh László (\url{nemethl@gyorsposta.hu})

